
%%%%%%%%%%%%%%%%%%%%%%%%%%%%%%%%%%%%%%%%%%%%%%%%%%%%%%%%%%%%%%%%%%%%%%%%%%%%%%%%%%%%%%%
%%%%%%%%%%%%%%%%%%%%%%%%%%%%%%%%%%%%%%%%%%%%%%%%%%%%%%%%%%%%%%%%%%%%%%%%%%%%%%%%%%%%%%%
% 
% This top part of the document is called the 'preamble'.  Modify it with caution!
%
% The real document starts below where it says 'The main document starts here'.

\documentclass[12pt]{article}

\usepackage{amssymb,amsmath,amsthm}
\usepackage[top=1in, bottom=1in, left=1.25in, right=1.25in]{geometry}
\usepackage{fancyhdr}
\usepackage{enumerate}
\usepackage{listings}
\usepackage{graphicx}
\usepackage{float}
\usepackage{multicol}
% Comment the following line to use TeX's default font of Computer Modern.
\usepackage{times,txfonts}
\usepackage{mwe}
\usepackage{caption}
\usepackage{subcaption}

\usepackage{tikz}
\def\checkmark{\tikz\fill[scale=0.4](0,.35) -- (.25,0) -- (1,.7) -- (.25,.15) -- cycle;} 



\makeatletter
\renewcommand*\env@matrix[1][*\c@MaxMatrixCols c]{%
  \hskip -\arraycolsep
  \let\@ifnextchar\new@ifnextchar
  \array{#1}}
\makeatother

\newtheoremstyle{homework}% name of the style to be used
  {18pt}% measure of space to leave above the theorem. E.g.: 3pt
  {12pt}% measure of space to leave below the theorem. E.g.: 3pt
  {}% name of font to use in the body of the theorem
  {}% measure of space to indent
  {\bfseries}% name of head font
  {:}% punctuation between head and body
  {2ex}% space after theorem head; " " = normal interword space
  {}% Manually specify head
\theoremstyle{homework} 

% Set up an Exercise environment and a Solution label.
\newtheorem*{exercisecore}{Exercise \@currentlabel}
\newenvironment{exercise}[1]
{\def\@currentlabel{#1}\exercisecore}
{\endexercisecore}

\newcommand{\localhead}[1]{\par\smallskip\noindent\textbf{#1}\nobreak\\}%
\newcommand\solution{\localhead{Solution:}}

%%%%%%%%%%%%%%%%%%%%%%%%%%%%%%%%%%%%%%%%%%%%%%%%%%%%%%%%%%%%%%%%%%%%%%%%
%
% Stuff for getting the name/document date/title across the header
\makeatletter
\RequirePackage{fancyhdr}
\pagestyle{fancy}
\fancyfoot[C]{\ifnum \value{page} > 1\relax\thepage\fi}
\fancyhead[L]{\ifx\@doclabel\@empty\else\@doclabel\fi}
\fancyhead[C]{\ifx\@docdate\@empty\else\@docdate\fi}
\fancyhead[R]{\ifx\@docauthor\@empty\else\@docauthor\fi}
\headheight 15pt

\def\doclabel#1{\gdef\@doclabel{#1}}
\doclabel{Use {\tt\textbackslash doclabel\{MY LABEL\}}.}
\def\docdate#1{\gdef\@docdate{#1}}
\docdate{Use {\tt\textbackslash docdate\{MY DATE\}}.}
\def\docauthor#1{\gdef\@docauthor{#1}}
\docauthor{Use {\tt\textbackslash docauthor\{MY NAME\}}.}
\makeatother

% Shortcuts for blackboard bold number sets (reals, integers, etc.)
\newcommand{\Reals}{\ensuremath{\mathbb R}}
\newcommand{\Nats}{\ensuremath{\mathbb N}}
\newcommand{\Ints}{\ensuremath{\mathbb Z}}
\newcommand{\Rats}{\ensuremath{\mathbb Q}}
\newcommand{\Cplx}{\ensuremath{\mathbb C}}
%% Some equivalents that some people may prefer.
\let\RR\Reals
\let\NN\Nats
\let\II\Ints
\let\CC\Cplx

%\textbf{Code:}
%\begin{center}
%  \lstinputlisting{NewtonsMethodP5.m}
%\end{center}
%
%\textbf{Console:}
%\begin{center}
%  \lstinputlisting{P5C.txt}
%\end{center}
%\vspace{.15in}


%\begin{figure}[H]
%  \begin{center}
%    \caption{The one-norm unit ball}
%    \includegraphics[width=.76\textwidth]{1norm.png}
%  \end{center}
%\end{figure}




%%%%%%%%%%%%%%%%%%%%%%%%%%%%%%%%%%%%%%%%%%%%%%%%%%%%%%%%%%%%%%%%%%%%%%%%%%%%%%%%%%%%%%%
%%%%%%%%%%%%%%%%%%%%%%%%%%%%%%%%%%%%%%%%%%%%%%%%%%%%%%%%%%%%%%%%%%%%%%%%%%%%%%%%%%%%%%%
% 
% The main document start here.

% The following commands set up the material that appears in the header.
\doclabel{Math 661: Homework 9}
\docauthor{Stefano Fochesatto}
\docdate{\today}


\begin{document}

\begin{exercise}{3.1} Apply symmetric rank-one quasi-Newton method to solve, 
  \begin{equation*}
    \text{minimize } f(x) = \frac{1}{2}x^TQx - c^Tx
  \end{equation*}
  with, 
  \begin{equation*}
    Q = \begin{pmatrix}
      5 & 2 & 1 \\
      2 & 7 & 3 \\
      1 & 3 & 9 
    \end{pmatrix}
    \text{ and }
    c = 
    \begin{pmatrix}
      -9\\
      0\\
      -8
    \end{pmatrix}
  \end{equation*}
  Initialize the method with $x_0 = (0, 0, 0)^T$ and $B_0 = I$. Use an exact line search. 
  \begin{proof}[Solution] 
    Applying symmetric rank one quasi-Newton method to solve this minimization problem begins with finding 
    a descent direction $p_k$ by solving the following system, 
    \begin{equation*}
      B_kp_k = - \nabla f(x_k).
    \end{equation*}
    Recall that in Homework \#7 we proved that the step size for exact line search on 
    a quadratic optimization problem is given by, 
    \begin{equation*}
      \alpha_k = \frac{-p_k^T \nabla f(x_k)}{p_k^TQp_k}.
    \end{equation*}
    We update $x$ using $\alpha_k$ from exact line search $x_{k + 1} = x_k + \alpha_k p_k$.
    Then we compute the following to update our estimated hermitian $B$,
    \begin{equation*}
      s_k = x_{k + 1} - x_k,
    \end{equation*}
    \begin{equation*}
      y_k = \nabla f(x_{k + 1}) - \nabla f(x_k),
    \end{equation*}
    \begin{equation*}
      B_{k + 1} = B_k + \frac{(y_k - B_ks_k)(y_k - B_ks_k)^T}{(y_k - B_ks_k)^Ts_k}.
    \end{equation*}
    Below is an implementation of the symmetric rank-one quasi-Newton method using exact line search applied to the given problem,\\ 
    \textbf{Code:}
    \begin{center}
      \lstinputlisting[basicstyle = \small]{sr1btquad.m}
    \end{center}
    
    \textbf{Console:}
    \begin{center}
      \lstinputlisting[basicstyle = \small]{r1.txt}
    \end{center}

  \end{proof}
\end{exercise}
\vspace{.5in}


\begin{exercise}{3.3} Let $f$ be a strictly convect quadratic function of one variable. Prove that the 
  secant method for minimization will terminate in exactly one iteration for any initial starting points
  $x_0$ and $x_1$. 
  \begin{proof} Suppose $f$ is a strictly convect quadratic function of one variable and let $x_0$ and $x_1$ be 
    initial points. Note that $f$ is of the form $f(x) = ax^2 + bx + c$ where $2a > 0$ and $f'$ is of the form 
    $f'(x) = 2ax + b$. Applying the secant method for minimization our first iteration looks like, 
    \begin{equation*}
      x_2 = x_1 - \frac{x_1 - x_0}{f'(x_1) - f'(x_0)}f'(x_1).
    \end{equation*}
    Checking if $f'(x_2) = 0$ we get, 
    \begin{align*}
      f'(x_2) &= 2a \left(x_1 - \frac{x_1 - x_0}{f'(x_1) - f'(x_0)}f'(x_1)\right) + b\\
      &= 2a \left(x_1 - \frac{x_1 - x_0}{2a(x_1) + b - (2a(x_0) + b)}(2a(x_1) + b)\right) + b\\
      &= 2a \left(x_1 - \frac{x_1 - x_0}{2a(x_1) - 2a(x_0)}(2a(x_1) + b)\right) + b\\
      &= 2a \left(x_1 - \frac{x_1 - x_0}{2a(x_1 - x_0)}(2a(x_1) + b)\right) + b\\
      &= 2a \left(x_1 - \left(\frac{2a(x_1)}{2a} + \frac{b}{2a}\right)\right) + b\\
      &= 2a \left(- \frac{b}{2a}\right) + b\\
      &= 0
    \end{align*}
    So the first order necessary conditions and the second order sufficient conditions are satisfied and we conclude that $x_2$ 
    is a local minimizer. 
  \end{proof}
\end{exercise}
\vspace{.5in}

\begin{exercise}{3.4} Let $C$ be a symmetric matrix of rank one. Prove that $C$ must have the form $C = \gamma w w^T$, where 
  $\gamma$ is a scalar and $w$ is a vector of norm one. 
  \begin{proof} Suppose $C$ is a symmetric matrix of rank one. Let $C = U\Sigma V^T$ be the singular value decomposition. Since 
    $C$ is symmetric we know that $C = C^T$ and therefore $U\Sigma V^T = V\Sigma U^T$. It then 
    follows that $U =  U\Sigma U^TV\Sigma^{-1} = U\Sigma V^TU\Sigma^{-1} = V$, and since $C$ is rank one $\Sigma$ has one non-zero 
    singular value, so our decomposition simplifies to $C = U\Sigma V^T = \sigma_1 uv^T$, and since $U = V$ we get $u = v$. 
  \end{proof}
\end{exercise}
\vspace{.5in}



\begin{exercise}{3.8} Let $B_{k + 1}$ be obtained from $B_k$ using the BFGS update formula. Prove that $B_{k + 1}s_k = y_k$. 
  \begin{proof} Recall that the BFGS update formula give $B_{k + 1}$ by, 
    \begin{equation*}
      B_{k + 1} = B_k - \frac{(B_ks_k)(B_ks_k)^T}{s_k^TB_ks_k} + \frac{y_ky_k^t}{y_k^Ts_k}.
    \end{equation*}
    Multiplying on the left by $s_k$ and simplifying we get the following, 
    \begin{align*}
      B_{k + 1}s_k &=  B_ks_k - \frac{(B_ks_k)(B_ks_k)^Ts_k}{s_k^TB_ks_k} + \frac{y_ky_k^T}{y_k^Ts_k}s_k,\\
     &=  B_ks_k - \frac{(B_ks_k)(s_k^TB_k^T)s_k}{s_k^TB_ks_k} + \frac{y_k(y_k^Ts_k)}{y_k^Ts_k},\\
     &=  B_ks_k - \frac{(B_ks_k)(s_k^TB_k^Ts_k)}{s_k^TB_ks_k} + \frac{y_k(y_k^Ts_k)}{y_k^Ts_k},\\
     &=  B_ks_k - B_ks_k + y_k,\\
     &=  y_k.
    \end{align*}

    \end{proof}
\end{exercise}
\vspace{.5in}


\begin{exercise}{2.1} Consider the problem, 
  \begin{align*}
    \text{minimize:  } &f(x) = x_1^2 + x_1^2x_3^2 + 2x_1x_2 + x_2^4 + 8x_2\\
    \text{subject to:  } &2x_1 + 5x_2 + x_3 = 3.
  \end{align*}
  \begin{enumerate}
    \item[(i)] Determine which of the following points are stationary points: $(0, 0, 2)^T$, $(0, 0, 3)^T$, and $(1, 0, 1)^T$. 
    \begin{proof}[Solution] Computing the gradient and Hessian, 
      \begin{equation*}
        \nabla f(x) = 
        \begin{pmatrix}
          2x_1 + 2x_1x_3^2 + 2x_2\\
          2x_1 + 4x_2^3 + 8\\
          2x_1^2x_3
        \end{pmatrix}
      \end{equation*}
      \begin{equation*}
        \nabla^2 f(x) = 
        \begin{pmatrix}
          2x_3^2 + 2 & 2 & 4x_1x_3\\
          2 & 12x_2^2 & 0\\
          4x_1x_3 & 0 & 2x_1^2
        \end{pmatrix}
      \end{equation*}
      Using the Lagrange Multipliers Necessary condition for optimality we need to see if there exists 
      a scalar $\lambda$ such that $\nabla f(x_*) = A^T\lambda$. Note that for our problem $A^T = (2, 5, 1)^T$. 
      Checking $(0, 0, 2)^T$ we get, 
      \begin{equation*}
        \nabla f(x) = 
        \begin{pmatrix}
          2(0) + 2(0)(2)^2 + 2(0)\\
          2(0) + 4(0)^3 + 8\\
          2(0)^2(0)
        \end{pmatrix} =         \begin{pmatrix}
          0\\
          8\\
          0
        \end{pmatrix}
        \neq
        \begin{pmatrix}
          2\\
          5\\
          1
        \end{pmatrix}
        \lambda.
      \end{equation*}
      Clearly there is no value for $\lambda$ which will make this expression true. So $(0, 0, 2)^T$
      is not a stationary point.

      Checking $(0, 0, 3)^T$ again we get,
      \begin{equation*}
        \nabla f(x) = 
        \begin{pmatrix}
          2(0) + 2(0)(2)^2 + 2(0)\\
          2(0) + 4(0)^3 + 8\\
          2(0)^2(0)
        \end{pmatrix} =         \begin{pmatrix}
          0\\
          8\\
          0
        \end{pmatrix}
        \neq
        \begin{pmatrix}
          2\\
          5\\
          1
        \end{pmatrix}
        \lambda.
      \end{equation*}
      Again we have the same situation where no value for $\lambda$ which will satisfy this expression so
      $(0, 0, 2)^T$ is not a stationary point. 
      
      Checking $(1, 0, 1)^T$ we get, 
      \begin{equation*}
        \nabla f(x) = 
        \begin{pmatrix}
          2(1) + 2(1)(1)^2 + 2(0)\\
          2(1) + 4(0)^3 + 8\\
          2(1)^2(1)
        \end{pmatrix} =         \begin{pmatrix}
          4\\
          10\\
          1
        \end{pmatrix}
         = 
        \begin{pmatrix}
          2\\
          5\\
          1
        \end{pmatrix}
        (2).
      \end{equation*}
      As we see when we set $\lambda = 2$ we get that expression is satisfied so $(1, 0, 1)^T$ is indeed a minimizer. 

    \end{proof}
    \vspace{.15in}

    \item[(ii)] Determine whether each stationary point is a local minimizer, local maximizer, or saddle point. 
    \begin{proof} To see what type of stationary point $x_* = (1, 0, 1)^T$ is we need to observe $Z^T\nabla^2f(x_*)Z$ where 
      $Z$ is the null space matrix for $A$. Note that $Z$ can look like the following, 
      \begin{equation*}
        Z = \begin{pmatrix}
          -\frac{1}{2} & -\frac{5}{2}\\
          0 & 1\\
          1 & 0
        \end{pmatrix}.
      \end{equation*}
      Computing $Z^T\nabla^2f(x_*)Z$ we get, 
      \begin{align*}
        Z^T\nabla^2f(x_*)Z & = \begin{pmatrix}
          -\frac{1}{2} & 0 & 1\\ 
          -\frac{5}{2} & 1 & 0
        \end{pmatrix}
        \begin{pmatrix}
          2(1)^2 + 2 & 2& 4(1)(1)\\
          2& 12(0)^2 & 0\\
          4(1)(1) & 0 & 2(1)^2
        \end{pmatrix}
        \begin{pmatrix}
          -\frac{1}{2} & -\frac{5}{2}\\
          0 & 1\\
          1 & 0
        \end{pmatrix}\\
        &=\begin{pmatrix}
          -\frac{1}{2} & 0 & 1\\ 
          -\frac{5}{2} & 1 & 0
        \end{pmatrix}
        \begin{pmatrix}
          4 & 2& 4\\
          2& 0 & 0\\
          4 & 0 & 2
        \end{pmatrix}
        \begin{pmatrix}
          -\frac{1}{2} & -\frac{5}{2}\\
          0 & 1\\
          1 & 0
        \end{pmatrix}\\
        &= 
        \begin{pmatrix}-1&-6\\ -6&15\end{pmatrix}
      \end{align*}
      Quickly solving for eigenvalues we get the following characteristic equation, 
      $(-1-\lambda)(15 - \lambda) - 36 = \lambda^2-14\lambda-51 = (\lambda + 3)(\lambda - 17)$
      which gives eigenvalues $-3$ and $17$ so the matrix is indefinite and $x_*$ is a saddle point. 
    \end{proof}
  \end{enumerate}
\end{exercise}
\vspace{.5in}


\begin{exercise}{P18} Apply the forward difference formula $f'(x) \approx (f(x + h) - f(x))/h$
  to estimate the gradient of the function, 
  \begin{equation*}
    f(x) = \exp(10x_1 + 2x_2^2)
  \end{equation*}
  at point $x = (-1, 1)$. Assuming that $\epsilon_{mach} = 2.2204 \times 10^{-16}$ on the computer 
  you used, how accurate is the approximated gradient when you actually compute using the best value of $h$ and the 'simpler'
  value of $h$, form page 425?
  \begin{proof}[Solution] Applying the feed forward difference formula we get an approximate gradient of, 
    \begin{equation*}
      \hat{\nabla f(x)} = 
      \begin{bmatrix}
        \frac{e^{10(x_1 + h) + 2x_2^2} - e^{10x_1 + 2x_2^2}}{h}\\
        \frac{e^{10x_1 + 2(x_2+h)^2} - e^{10x_1 + 2x_2^2}}{h}
      \end{bmatrix}
    \end{equation*}
    Computing our estimated gradient with $hs = \sqrt{\epsilon_{mach}}$(simple h) value we get the following, \\

      \textbf{Console:}
      \begin{center}
        \lstinputlisting[basicstyle = \small]{r2.txt}
      \end{center}
    The 'complex' value for $h$ is given by the following, 
    \begin{equation*}
      h_1 = \sqrt{\frac{2|f(x)|\epsilon}{|\frac{d^2f}{dx_1^2}(x)|}},
    \end{equation*}
    \begin{equation*}
      h_2 = \sqrt{\frac{2|f(x)|\epsilon}{|\frac{d^2f}{dx_2^2}(x)|}}.
    \end{equation*}
    The partial derivatives that we will use are, 
    \begin{equation*}
      \frac{d^2f}{dx_1^2}(x) = 100e^{10x_1 + 2x_2^2},
    \end{equation*}
    \begin{equation*}
      \frac{d^2f}{dx_2^2}(x) = e^{10x_1 + 2x_2^2}(16x_2^2 + 4).
    \end{equation*}
    Computing the separate $h$ values and plugging into our feed forward difference formula we get the following,\\

    \textbf{Console:}
    \begin{center}
      \lstinputlisting[basicstyle = \small]{r3.txt}
    \end{center}

    Comparing our two approximation of the gradient we find that the more 'complex' $h$ value led to a single order of magnitude smaller 
    error than the 'simple' $h$ value.\\


    \textbf{Console:}
    \begin{center}
      \lstinputlisting[basicstyle = \small]{r4.txt}
    \end{center}
  \end{proof}
\end{exercise}
\vspace{.5in}



\begin{exercise}{P19} Suppose a user want to solve $min_{x \in \RR}$ for a smooth objective function $f: \RR^n \to \RR$, 
  but they do not provide a gradient function $\nabla f(x)$. We can still use quasi-Newton if we apply finite differences to replace
  the missing gradient. Use this idea to write a new code, 
  \begin{equation*}
    \text{function [xk, xklist] = sr1btfd(x0, f, tol, maxiters)}
  \end{equation*}
  which replaces the user-provided gradient in sr1bt.m with a finite-difference gradient. Test it on the Rosenbrock example, 
  as is done rosencompare.m, and give iterations. (Don't worry about visualizing, or the contours.) How would it help if the user 
  provided values for the typical scales of the gradient and the Hessian. 
  \begin{proof}[Solution] The following is a modification of the sr1bt.m code removing instances where the gradient is computed with the 
    user supplied function, now we estimate the gradient with a finite difference gradient using a simple value for $h = \sqrt{\epsilon_{mach}}$.\\
    
    
    \textbf{Code:}
    \begin{center}
      \lstinputlisting[basicstyle = \small]{sr1btfd.m}
    \end{center}

    Testing this code on the Rosenbrock example we get 26 iterations, just like if we had provided the gradient information 
    as in SR1BT. (They actually do also provide different iterations, I checked by subtracting the xklist objects and seeing if they are nonzero) \\

       
    \textbf{Code:}
    \begin{center}
      \lstinputlisting[basicstyle = \small]{r5.txt}
    \end{center}
    
    If the user were to provide information about the scales fo the gradient and the Hessian, we could see if $\sqrt{\epsilon_{mach}}$
    is a good approximation in the finite differences formula for $h$. Recall that in the finite differences formula 
    we have a more complex $h$ value that is given by, 
    \begin{equation*}
      h_i = \sqrt{\frac{2|f(x)|\epsilon_{mach}}{\frac{d^2f}{dx_i^2}(x)}}.
    \end{equation*}
    The scale of the Hessian, or it's condition number would tell us how much these values for $h$ vary in the finite difference formula. 
    A condition number close to 1 suggests that all second order partial derivatives are of the same order of magnitude and therefore 
    $h_i$ are relatively similar in terms of size, in which case we might use the 'simple' formula for $h$. More directly if the value of the 
    function and the value of the diagonal elements of the Hessian are of the same scale then the 'simple' $h$ will be suitable. 
  \end{proof}
\end{exercise}




\end{document}




























